_{t+1}=\alpha g_t$, this implies
\[
    d^\top g_t
    \leq
    -
    \frac{1}{\alpha L}
    \|d\|_2^2.
\]
Consequently,
\begin{align}
    \|g_{t+1}\|_2^2
    &=
    \|g_t+d\|_2^2 \\
    &=
    \|g_t\|_2^2
    +
    2d^\top g_t
    +
    \|d\|_2^2 \\
    &\leq
    \|g_t\|_2^2
    +
    \left(
        1-\frac{2}{\alpha L}
    \right)
    \|d\|_2^2.
\end{align}
Because $\alpha L\leq1$, the final coefficient is nonpositive, and hence
\[
    R_{t+1}
    =
    \|g_{t+1}\|_2^2
    \leq
    \|g_t\|_2^2
    =
    R_t.
\]
This proves both claims.
\end{proof}

\subsection{Quantitative Gaussian Reuse Rates}
\label{app:gaussian_reuse_rate}

\begin{proof}[Proof of Proposition~\ref{prop:reuse_rate}]
Let \(\delta_t=\mu_t-a\) and
\(X_t=\|\delta_t\|_2^2/(2\sigma^2)=D_t-C_\Sigma\).
For an uncontrolled negative sample of mass \(q\),
\[
\delta_{t+1}
=\left(1+\frac{\eta q}{\sigma^2}\right)\delta_t,
\]
which gives the stated exact geometric recurrence.

Now let \(\omega_t=e^{-a[D_t-\tau]_+}\). Once
\(D_t>\tau\), write
\(b=\eta q/\sigma^2\) and
\(K_0=e^{a(\tau-C_\Sigma)}\). Then
\begin{equation}
X_{t+1}=X_t\left(1+bK_0e^{-aX_t}\right)^2.
\label{eq:app_tapered_gaussian_recurrence}
\end{equation}
Consequently \(\Delta X_t\ge0\), and for all sufficiently large
\(X_t\),
\begin{equation}
0\le\Delta X_t\le C X_te^{-aX_t}
\label{eq:app_gaussian_increment_bound}
\end{equation}
for a finite constant \(C\). Put \(Y_t=e^{aX_t}\). Since
\(aX_te^{-aX_t}\le1/e\),
Equation~\eqref{eq:app_gaussian_increment_bound} also gives a uniform
bound \(a\Delta X_t\le M\). Hence
\begin{align}
\Delta Y_t
&=Y_t(e^{a\Delta X_t}-1) \\
&\le C_1Y_t\,a\Delta X_t
\le C_2\log Y_t.
\label{eq:app_gaussian_y_increment}
\end{align}
The standard comparison sequence
\(\overline Y_t=K t\log t\), with \(K\) chosen larger than the
constant in Equation~\eqref{eq:app_gaussian_y_increment}, is eventually
a supersolution. Therefore
\(Y_t=O(t\log t)\) and
\[
X_t\le a^{-1}\{\log t+\log\log t+O(1)\}.
\]
If an application does not make the sign of \(\Delta X_t\) explicit,
negative increments already satisfy the desired upper comparison, while
the nonnegative branch obeys the argument above.

In continuous time, the far-branch equation is
\(\dot X=C_3Xe^{-aX}\). Thus
\(Y=e^{aX}\) satisfies \(\dot Y=C_3\log Y\), and the same comparison
gives \(X(t)=O(a^{-1}\log t)\).
\end{proof}

\subsection{Far-Field Radial Bridge}
\label{app:far_field_bridge}

\begin{proof}[Proofs of Lemma~\ref{lem:far_field_radial_balance} and
Corollary~\ref{cor:selective_vs_global}]
Set \(\mu=m_++rv\) with \(\|v\|_2=1\). Direct expansion gives
\begin{align}
\frac{\langle v,F(\mu)\rangle}{r}
={}&-p\,v^\top\Sigma^{-1}v
+\sum_iq_i v^\top\Sigma^{-1}v \\
&+\frac1r\sum_iq_i
\langle v,\Sigma^{-1}(m_+-a_i)\rangle.
\end{align}
The final term vanishes, proving the lemma.

With atom-specific weights, \(D_\mu(a_i)\to\infty\) along every ray.
For finitely many atoms, \(\omega_i(D_i)\to0\) therefore makes both the
weighted coefficient of \(r\) and the weighted constant term vanish.
The radial limit is
\(-p v^\top\Sigma^{-1}v\), which is strictly negative. This is a
pointwise-in-\(q/p\) statement: the radius at which it becomes negative
can grow with \(q/p\). Moreover,
\(\|\Sigma^{-1}(\mu-a_i)\|_2=O(\sqrt{D_i})\), so
\(\omega(D)\sqrt D\to0\) removes the absolute negative force. Finally,
for \(\omega_i\equiv\alpha\), the lemma applies with \(q\) replaced by
\(\alpha q\), yielding the condition \(\alpha q<p\).
\end{proof}

\subsection{Reference-Regularized Restoration}
\label{app:regularized_restoration}

\begin{proof}[Proof of Proposition~\ref{prop:regularized_restoration}]
The unweighted Gaussian field with a KL-geometry spring is
\[
F_\gamma(\mu)=\Sigma^{-1}
\{(q-p-\gamma)\mu+pm_+-qm_-+\gamma\mu_0\}.
\]
Whenever \(p-q+\gamma\ne0\), its unique equilibrium is
\[
\mu_\gamma^*
=\frac{pm_+-qm_-+\gamma\mu_0}{p-q+\gamma},
\]
and the field Jacobian is
\((q-p-\gamma)\Sigma^{-1}\). It is Hurwitz exactly when
\(\gamma>q-p\). This yields all three regimes in the proposition; for
\(p=q\), the displayed equilibrium has displacement proportional to
\(1/\gamma\). When \(p>q\), increasing \(\gamma\) contracts the
equilibrium toward \(\mu_0\), potentially sacrificing extrapolation.

For the Euclidean spring \(-\gamma(\mu-\mu_0)\), the Jacobian is
\((q-p)\Sigma^{-1}-\gamma I\). In the negative-dominant regime it is
Hurwitz exactly when
\(\gamma>(q-p)\lambda_{\max}(\Sigma^{-1})\); this condition is not
interchangeable with the KL-geometry coefficient above.
\end{proof}

\subsection{Gaussian Closed-Loop Boundedness}
\label{app:gaussian_drpo_boundedness}

\begin{proof}[Proof of Theorem~\ref{thm:gaussian_drpo_boundedness}]
Let \(e=\mu-m_+\) and
\[
N_\omega(\mu)=\sum_iq_i\omega_i(D_i)
\Sigma^{-1}(\mu-a_i).
\]
The finite-atom exponential-taper field is locally Lipschitz and has at
most linear growth, so its solutions exist for all forward times.
For \(V(e)=\frac12\|e\|_2^2\),
\begin{align}
\dot V
&=-p e^\top\Sigma^{-1}e+e^\top N_\omega(\mu) \\
&\le-p\lambda_{\min}(\Sigma^{-1})\|e\|_2^2
+B_\omega\|e\|_2.
\label{eq:app_gaussian_lyapunov_bound}
\end{align}
Thus \(\dot V<0\) outside every ball of radius
\(R^*+\varepsilon\). The standard ultimate-boundedness argument gives
\[
\limsup_{t\to\infty}\|\mu(t)-m_+\|_2\le R^*.
\]
On the boundary of any ball centered at \(m_+\) with radius greater
than \(R^*\), Equation~\eqref{eq:app_gaussian_lyapunov_bound} is
strictly negative. Continuity of the field and the Poincar\'e--Bohl
inward-pointing theorem therefore imply at least one zero inside the
ball. This is an existence result, not a uniqueness or global-convergence
claim.

It remains to verify \(B_\omega<\infty\) and the stated constant. Put
\(r_i=\|\Sigma^{-1/2}(\mu-a_i)\|_2\), so
\(D_i=C_\Sigma+r_i^2/2\), and
\[
\|\Sigma^{-1}(\mu-a_i)\|_2
\le\sqrt{\lambda_{\max}(\Sigma^{-1})}\,r_i.
\]
For \(a=\lambda/c\) and
\(\widetilde\tau_i=\tau_i-C_\Sigma\), maximize
\[
g_i(r)=r\exp\{-a[r^2/2-\widetilde\tau_i]_+\},\qquad r\ge0.
\]
If \(\widetilde\tau_i<0\), the taper is active at zero and the maximum
occurs at \(r=a^{-1/2}\), giving
\(a^{-1/2}e^{a\widetilde\tau_i-1/2}\). If
\(\widetilde\tau_i\ge0\), the untapered plateau ends at
\(r_\tau=\sqrt{2\widetilde\tau_i}\). When
\(a^{-1/2}\le r_\tau\), the maximum is \(r_\tau\); otherwise it is the
same interior-tail value. These are exactly the three displayed
branches for \(G_i\), and the triangle inequality proves the bound on
\(B_\omega\).

For near-field fidelity, let \(\mu^*\) be an uncontrolled equilibrium
with \(D_{\mu^*}(a_i)<\tau_i\) for every negative atom. Strictness and
continuity provide a neighborhood in which every \(\omega_i=1\).
The controlled and uncontrolled fields, including their Jacobians, are
therefore identical there.

Finally, along any sequence \(\|\mu_n\|_2\to\infty\), each fixed atom
satisfies \(D_{\mu_n}(a_i)\to\infty\), hence
\(\omega_i(D_{\mu_n}(a_i))\to0\). Finiteness of the atom set gives
\(\rho_{\mathrm{eff}}(\mu_n)\to0\).
\end{proof}

\subsection{Derivation and Proof of
Theorem~\ref{thm:aggregate_equilibria}}
\label{app:aggregate_equilibrium}

We use a regular minimal exponential-family representation to cover the
continuous-action Gaussian and discrete-action categorical policies in a
common set of output coordinates. At a fixed context, write
\begin{equation}
    \pi_\xi(a)
    =
    h(a)
    \exp\left\{
        \xi^\top T(a)-\psi(\xi)
    \right\},
    \qquad
    \xi\in\mathcal H,
    \label{eq:app_exponential_family_policy}
\end{equation}
where \(\xi\) is the natural parameter, \(T(a)\) is the sufficient
statistic, and \(\psi\) is the log-partition function. Let
\begin{equation}
    \mathcal M
    =
    \left\{
        \nabla\psi(\xi):
        \xi\in\mathcal H
    \right\}
    \label{eq:app_mean_parameter_space}
\end{equation}
denote the attainable mean-parameter space.

For a fixed-covariance Gaussian mean policy,
\(\xi=\Sigma^{-1}\mu\), \(T(a)=a\), and
\(\nabla\psi(\xi)=\mu\). For a categorical policy, fixing a
reference-logit gauge gives a minimal \((K-1)\)-dimensional
representation in which \(T(a)\) is the corresponding reduced action
indicator and \(\nabla\psi(\xi)\) is the reduced probability vector.
Thus, the two cases differ in their sufficient statistics and attainable
mean spaces, but share the derivation below.

Write
\[
    \widehat A
    =
    \widehat A^+
    -
    \widehat A^-,
    \qquad
    \widehat A^+
    =
    \max\{\widehat A,0\},
    \qquad
    \widehat A^-
    =
    \max\{-\widehat A,0\}.
\]
Define the aggregate positive and negative update masses and their
unnormalized moments:
\begin{equation}
\begin{aligned}
    p
    &=
    \mathbb E_\nu[\widehat A^+(a)],
    \qquad
    q
    =
    \mathbb E_\nu[\widehat A^-(a)],
    \\
    \bar{\mathbf m}_+
    &=
    \mathbb E_\nu
    \left[
        \widehat A^+(a)T(a)
    \right],
    \\
    \bar{\mathbf m}_-
    &=
    \mathbb E_\nu
    \left[
        \widehat A^-(a)T(a)
    \right].
\end{aligned}
\label{eq:app_aggregate_statistics}
\end{equation}
When \(p>0\), define
\begin{equation}
    \mathbf m_+
    =
    \frac{\bar{\mathbf m}_+}{p}.
    \label{eq:app_positive_moment}
\end{equation}
This is the positive-only target used in the main text: it is the
mean-parameter point selected by positive feedback alone, provided that
it lies in \(\mathcal M\). When \(q>0\), define
\begin{equation}
    \mathbf m_-
    =
    \frac{\bar{\mathbf m}_-}{q}.
    \label{eq:app_negative_moment}
\end{equation}
This is the negative-feedback moment, namely the advantage-weighted
center of the negative feedback. It is not an attraction target; the
negative actor contribution repels the learner from this moment. The
unnormalized quantities
\(\bar{\mathbf m}_+\) and \(\bar{\mathbf m}_-\) remain well defined in
the boundary cases \(p=0\) or \(q=0\).

The signed population objective is
\begin{equation}
    \mathcal L(\xi)
    =
    \mathbb E_\nu
    \left[
        \widehat A(a)
        \log\pi_\xi(a)
    \right].
    \label{eq:app_signed_population_objective}
\end{equation}
Using Equation~\eqref{eq:app_exponential_family_policy} and
\(\widehat A=\widehat A^+-\widehat A^-\), we obtain
\begin{equation}
    \mathcal L(\xi)
    =
    \left(
        \bar{\mathbf m}_+
        -
        \bar{\mathbf m}_-
    \right)^\top\xi
    -
    (p-q)\psi(\xi)
    +
    C,
    \label{eq:app_signed_objective}
\end{equation}
where
\[
    C
    =
    \mathbb E_\nu
    \left[
        \widehat A(a)\log h(a)
    \right]
\]
is independent of \(\xi\). Therefore, the aggregate actor field is
\begin{equation}
    \mathbf F(\xi)
    =
    \nabla\mathcal L(\xi)
    =
    \bar{\mathbf m}_+
    -
    \bar{\mathbf m}_-
    -
    (p-q)\nabla\psi(\xi),
    \label{eq:app_aggregate_field}
\end{equation}
and
\begin{equation}
    \nabla^2\mathcal L(\xi)
    =
    -(p-q)\nabla^2\psi(\xi).
    \label{eq:app_signed_hessian}
\end{equation}
Regularity and minimality imply that
\(\nabla^2\psi(\xi)\succ0\) throughout \(\mathcal H\), and that
\(\nabla\psi\) is one-to-one from \(\mathcal H\) onto
\(\mathcal M\).

We analyze the continuous dynamics and discrete update
\begin{equation}
    \dot\xi
    =
    \mathbf F(\xi),
    \qquad
    \xi_{t+1}
    =
    \xi_t
    +
    \alpha\mathbf F(\xi_t),
    \qquad
    \alpha>0.
    \label{eq:app_aggregate_dynamics}
\end{equation}

\begin{proof}

\medskip
\noindent\textbf{Positive-only case.}
Suppose \(p>0\) and \(q=0\). Then
\(\bar{\mathbf m}_-=\mathbf 0\), and
Equation~\eqref{eq:app_aggregate_field} becomes
\begin{equation}
    \mathbf F(\xi)
    =
    p
    \left(
        \mathbf m_+
        -
        \nabla\psi(\xi)
    \right).
    \label{eq:app_positive_only_field}
\end{equation}
A finite equilibrium must therefore satisfy
\begin{equation}
    \nabla\psi(\xi^\star)
    =
    \mathbf m_+.
    \label{eq:app_positive_only_equilibrium}
\end{equation}
If \(\mathbf m_+\in\mathcal M\), existence follows because
\(\nabla\psi\) maps \(\mathcal H\) onto \(\mathcal M\), and uniqueness
follows because this map is one-to-one.

The continuous field Jacobian is
\begin{equation}
    J_{\mathbf F}(\xi^\star)
    =
    -p\nabla^2\psi(\xi^\star)
    \prec0,
    \label{eq:app_positive_only_continuous_jacobian}
\end{equation}
so the equilibrium is locally asymptotically stable. For the discrete
update, the local Jacobian is
\begin{equation}
    J_{\mathrm{disc}}(\xi^\star)
    =
    I
    -
    \alpha p\nabla^2\psi(\xi^\star).
    \label{eq:app_positive_only_discrete_jacobian}
\end{equation}
Its eigenvalues have magnitude smaller than one whenever
\begin{equation}
    0
    <
    \alpha
    <
    \frac{
        2
    }{
        p\lambda_{\max}
        \left(
            \nabla^2\psi(\xi^\star)
        \right)
    }.
    \label{eq:app_positive_only_stepsize}
\end{equation}
Hence, the positive-only equilibrium is locally stable under
sufficiently small discrete steps. If
\(\mathbf m_+\notin\mathcal M\), no finite equilibrium exists.

\medskip
\noindent\textbf{Positive-dominant regime.}
Suppose \(p>q>0\). Since both normalized moments are defined, a finite
equilibrium must satisfy
\begin{equation}
    \mathbf F(\xi^\star)
    =
    \mathbf 0,
\end{equation}
or equivalently
\begin{equation}
\begin{aligned}
    \nabla\psi(\xi^\star)
    &=
    \frac{
        \bar{\mathbf m}_+
        -
        \bar{\mathbf m}_-
    }{
        p-q
    }
    \\
    &=
    \frac{
        p\mathbf m_+
        -
        q\mathbf m_-
    }{
        p-q
    }
    =
    \mathbf m^\star.
\end{aligned}
\label{eq:app_equilibrium_condition}
\end{equation}
If \(\mathbf m^\star\in\mathcal M\), existence and uniqueness again
follow from the bijectivity of
\(\nabla\psi:\mathcal H\rightarrow\mathcal M\).

Since \(p-q>0\), Equation~\eqref{eq:app_signed_hessian} gives
\begin{equation}
    \nabla^2\mathcal L(\xi)
    \prec0.
\end{equation}
Thus, \(\mathcal L\) is strictly concave, and
\(\xi^\star\) is its unique finite maximizer.

For the continuous dynamics, the field Jacobian at the equilibrium is
\begin{equation}
    J_{\mathbf F}(\xi^\star)
    =
    -(p-q)
    \nabla^2\psi(\xi^\star)
    \prec0.
    \label{eq:app_continuous_jacobian}
\end{equation}
All eigenvalues are strictly negative, so
\(\xi^\star\) is locally asymptotically stable.

For the discrete update, the local Jacobian is
\begin{equation}
    J_{\mathrm{disc}}(\xi^\star)
    =
    I
    -
    \alpha(p-q)
    \nabla^2\psi(\xi^\star).
    \label{eq:app_discrete_jacobian}
\end{equation}
Let \(\lambda_i>0\) denote the eigenvalues of
\(\nabla^2\psi(\xi^\star)\). The corresponding eigenvalues of
\(J_{\mathrm{disc}}(\xi^\star)\) are
\(1-\alpha(p-q)\lambda_i\). Their magnitudes are strictly smaller than
one whenever
\begin{equation}
    0
    <
    \alpha
    <
    \frac{
        2
    }{
        (p-q)
        \lambda_{\max}
        \left(
            \nabla^2\psi(\xi^\star)
        \right)
    }.
    \label{eq:app_discrete_stepsize}
\end{equation}
This proves local discrete-time stability for sufficiently small
\(\alpha\).

Define \(\rho=q/p\). Algebraically,
\be